\chapter{El contrato social}
\label{ch:welfare}

\section{El compromiso}

Cuba conserva un Estado de bienestar universal al nivel de Europa occidental
---salud, educación, pensiones, discapacidad, apoyo a la infancia, sostén de
ingresos--- y el resto de la Parte III se ordena en torno a poder pagarlo.

Ese es el único elemento no negociable del libro, y este capítulo es donde se
argumenta en vez de afirmarse. Cuatro razones, en orden descendente de peso.

\emph{Es lo que vale la pena conservar.} El capítulo~\ref{ch:ledgerch} auditó
sesenta y siete años y encontró un activo inequívoco. Los resultados de salud están
medidos, no proclamados; los de educación los evaluó un organismo externo y superan
a la región en dos desviaciones estándar, con el cuartil más pobre por encima del
promedio regional. Nada más en ese balance está en esa categoría.

\emph{Es lo que hace políticamente posible el resto.} Una población a la que se le
pide aceptar mercados, propietarios extranjeros, liberalización de precios y
privatización los aceptará si el piso aguanta. No los aceptará si el piso
desaparece, y la evidencia no es teórica: las transiciones poscomunistas que
retiraron temprano la protección social vieron sus reformas económicas revertidas
por sus propios electorados en una década \citep{roland2000}, que es la manera más
cara de descubrir que el consentimiento es un insumo. A Cuba se le está pidiendo
correr una transición con una población a la que se le ha dicho durante sesenta y
siete años que esto es lo que compró el sacrificio. Quítelo en el año dos y no hay
año cinco.

\emph{Es la protección de aquellos a quienes la transición va a hacer daño.} El
capítulo~\ref{ch:premises} identificó a los expuestos: hogares sin remesas, que son
desproporcionadamente negros; las provincias orientales; los viejos; el campo. Toda
medida liberalizadora de este libro favorece al hogar con capital y con un pariente
en el exterior. El piso universal es el instrumento que impide que eso se convierta
en una estratificación permanente, y es el único que funciona sin una capacidad
administrativa que Cuba no tiene.

\emph{Es un insumo, no solo una transferencia.} El sistema educativo es aquello que
los capítulos~\ref{ch:talent} y \ref{ch:life} proponen monetizar. No puede
desfinanciarse una década y reiniciarse, porque los maestros se van y las cohortes
pasan. Gastar en él no es consumir el producto del modelo de crecimiento; es la
materia prima del modelo de crecimiento.

\section{Lo que cambia es quién paga}

Aquí está el movimiento central del capítulo, y la frase que conecta la Parte III con
las Partes I y II.

\emph{El Estado de bienestar cubano nunca lo han financiado los cubanos.}

Lo financiaron entre 1972 y 1989 los términos de intercambio soviéticos ---azúcar
comprado por encima del precio mundial, petróleo suministrado por debajo, y una
renta petrolera que a Cuba se le permitía reexportar---. Lo financiaron entre 2004 y
2015 aproximadamente el petróleo venezolano y los contratos de servicios médicos que
esa relación petrolera generó. En los intervalos lo financiaron la emisión monetaria
y el no mantener nada, que es la razón de que los hospitales no tengan agua y de que
La Habana pierda edificios por derrumbe.

Una prestación financiada por una renta externa dura exactamente lo que dura la
renta. Cuba ya lo ha descubierto dos veces, en 1990 y otra vez después de 2015, y en
ambas ocasiones el Estado de bienestar se salvó solo con un esfuerzo extraordinario y
a costa de su planta física. Los sistemas europeos que este libro toma como norma se
financian gravando una base doméstica ---que es por lo que sobreviven a los cambios
de gobierno, de precio de las materias primas y de patrón, y por lo que ningún
gobierno europeo ha tenido nunca que elegir entre pagar pensiones e importar
combustible---.

La propuesta es entonces: la misma prestación, otra fontanería. Una estructura
europea de prestaciones sobre una estructura europea de impuestos. Lo financiado por
renta pasa a financiarse por impuestos.

No es una empresa menor de lo que suena. Cuba no tiene un sistema tributario moderno
---abolió la mayor parte de la tributación en los años sesenta con el argumento de
que un Estado dueño de todo no necesita gravarse a sí mismo, y reintrodujo impuestos
solo a partir de 1994 y de forma parcial---. No tiene cultura de declaración, ni
historial significativo de fiscalización, ni infraestructura de retención fuera de la
nómina estatal, ni una población que haya pagado impuesto sobre la renta en memoria
viva. Construir una administración tributaria capaz de cobrar un impuesto al valor
agregado ha llevado de cinco a ocho años en todas partes donde se ha intentado. Es el
plazo de maduración más largo de este libro y por eso el
capítulo~\ref{ch:sequence} lo arranca en la primera fase, cuando no rinde nada.

\section{La aritmética}

Las cifras que siguen son órdenes de magnitud con sus supuestos declarados, no
pronósticos. Están aquí porque un compromiso de este tipo carece de sentido sin una
respuesta aproximada a cuánto cuesta y a si puede pagarse.

\emph{Cuánto cuesta la prestación.} Un sistema universal de nivel europeo se sitúa en
torno al 25 o 30 por ciento del ingreso nacional en gasto social: salud pública en
torno al 7 u 8 por ciento, educación en torno al 5, pensiones entre el 10 y el 12, y
el resto de protección social el saldo. El gasto público total en los Estados
miembros de la Unión Europea se sitúa en la franja del 40 al 50 por ciento y los
ingresos tributarios entre el 35 y el 45 por ciento.\unv{}

\emph{Qué hace Cuba hoy.} Las cifras publicadas cubanas sitúan el gasto social en una
proporción muy alta del producto, y los ingresos también, pero ---según la nota sobre
las fuentes--- la mayor parte de eso es un artefacto de convertir una economía de
propiedad estatal a un tipo de cambio administrativo, y esas proporciones no
significan lo que significan las mismas proporciones en otros lugares. Lo que sí puede
decirse sin el problema de la conversión es físico: Cuba entrega hoy algo cercano a la
cobertura universal a un costo real de recursos muy inferior al europeo, porque paga a
sus médicos, enfermeros y maestros una fracción de un salario europeo y ha dejado de
mantener sus edificios y de reponer su equipamiento. El sistema está subsidiado por su
propio personal y por su propio acervo de capital.

\emph{Dónde está la brecha.} Esa es la forma honesta del problema. Cuba no necesita
comprar cobertura europea ---tiene cobertura---. Necesita comprar \emph{insumos}
europeos: salarios que impidan que se vaya el personal, medicinas que estén de verdad
en la farmacia, equipos que funcionen, edificios que no se inunden, y una pensión con
la que alguien pueda vivir. Esas son las partes caras y son exactamente las que hoy
están sin financiar.

\emph{Si es asequible ahora.} No. Con el producto cubano de 2026 un nivel europeo de
prestaciones no es financiable a ninguna tasa impositiva, y un libro que fingiera otra
cosa no valdría nada. Lo que se propone en su lugar es una trayectoria declarada:
proteger el \emph{piso} de inmediato y dejar que el \emph{techo} converja a medida que
crezca la base, con la convergencia dotada de una meta y una fecha en vez de dejarla en
aspiración, y con el sistema tributario construido ahora para que pueda recaudar
cuando la base exista. El diseño que sigue está escrito para ser asequible en su
primera fase y para escalar.

\section{Qué modelo europeo}

Tres familias, y la elección importa más de lo que parece.

El modelo \emph{nórdico} es universal, financiado en gran medida con ingresos
generales y un impuesto al consumo amplio y alto, intensivo en servicios más que en
transferencias monetarias, y depende de una tasa de empleo alta. El modelo
\emph{bismarckiano} de Alemania y Francia es contributivo y se financia con nóminas,
de modo que la cobertura sigue al empleo formal. El patrón \emph{sureuropeo} es
universal en la ley y desigual en la prestación, financiado por una base estrecha
junto a una gran informalidad.

Cuba derivaría por defecto hacia el tercero. Es por tanto el resultado contra el cual
diseñar, y las decisiones que siguen se toman específicamente para evitarlo.

El libro propone una \emph{estructura nórdica de prestaciones}: universal, intensiva
en servicios, sin comprobación de medios, y financiada por una \emph{base amplia de
consumo y propiedad} antes que por nóminas.

Universalidad en vez de comprobación de medios, por dos razones que son sobre Cuba y
no sobre principios. Cuba ya tiene universalidad y desmantelarla para construir un
aparato de focalización destruiría un sistema que funciona para levantar uno que no
tiene capacidad de operar. Y toda comprobación de medios es una decisión discrecional
que toma un funcionario sobre un ciudadano, que es precisamente la superficie que el
capítulo~\ref{ch:integrity} dedica su extensión a eliminar. En un país donde
\emph{resolver} es un verbo, una prestación que depende de la evaluación de un
funcionario es una prestación que depende de conocer al funcionario.

En contra del financiamiento por nóminas como instrumento principal, dos razones más.
La cobertura seguiría al empleo formal en una economía donde la mayoría de los empleos
nuevos empezarán informales, lo que reproduce directamente el resultado sureuropeo. Y
una cuña alta sobre la nómina es el modo más fiable conocido de matar la pequeña
empresa, que es el sector que el capítulo~\ref{ch:capital} necesita hacer crecer y que
Cuba lleva cincuenta años sin tener.

\begin{design}{La estructura de prestaciones}
\label{d:benefit}
\begin{rules}
\rl{1}{Universal en especie} Atención de salud y educación gratuitas en el punto de
uso para todo residente, sin comprobación de medios, sin condición contributiva ni
requisito de nacionalidad más allá de la residencia. Este es el arreglo cubano
existente y se conserva sin cambio.

\rl{2}{Universal en efectivo} Una pensión básica plana universal a un nivel fijado
para prevenir la pobreza, pagadera a todo residente por encima de la edad de retiro
con independencia del historial contributivo; una prestación universal por hijo; y
apoyo por discapacidad según necesidad evaluada y no según contribución.

\rl{3}{Sustitución de ingresos por separado} Por encima del tramo plano, una pensión
ligada a los ingresos en cuentas individuales, capitalizadas o nocionales, con la
contribución visible en la nómina. Separar los dos tramos es lo que permite que el
tramo plano sea suficiente sin ser inasequible.

\rl{4}{Desempleo} Un seguro de desempleo con límite temporal y componente activo,
introducido en la fase en que se reduzca el empleo estatal y no antes ---porque es el
instrumento que vuelve políticamente posible esa reducción, e introducirlo tarde es lo
que convirtió el plan de plantilla de 2011 en un plan que nadie ejecutó---.

\rl{5}{Legisladas, no administradas} Todo nivel de prestación se fija por ley y se
actualiza por fórmula publicada. Ninguna prestación es variable a discreción
administrativa. Esa es la diferencia entre un derecho y un favor, y según el
Diseño~\ref{d:commit} es además un dispositivo de compromiso.
\end{rules}
\end{design}

\section{El sistema tributario}

Este es el núcleo operativo del capítulo. El capítulo~\ref{ch:money} es dueño de la
estabilización, la moneda y el déficit; este capítulo es dueño de qué se grava y a
qué tasa. Los dos tienen que coincidir en una cifra ---la proporción de ingresos que
exige el compromiso de bienestar--- y el capítulo~\ref{ch:money} es donde se muestra
que es financiable sin inflación.

El principio de diseño en todo momento es que una administración débil recauda un
impuesto sencillo y fracasa con uno ingenioso. Cuba está construyendo una
administración tributaria casi desde cero, en una economía de efectivo, con un
funcionariado desmoralizado y sin cultura de declaración. Cada sofisticación es un
lugar por donde el sistema pierde o donde un funcionario adquiere discrecionalidad.

\begin{design}{La estructura de ingresos}
\label{d:tax}
\begin{rules}
\rl{1}{Base amplia, tasas bajas, sin exenciones} Se enuncia primero porque gobierna
todo lo demás. Las exenciones son el modo en que mueren los sistemas tributarios: cada
una es pequeña, cada una es defendible, cada una crea una frontera sobre la que
discutir, y la discusión se conduce entre un contribuyente y un funcionario con
discrecionalidad. Sin exoneraciones sectoriales, sin tasas negociadas, sin regímenes
especiales para inversionistas nominados.

\rl{2}{El IVA como pieza central} Un impuesto al valor agregado de tasa única, una
sola tasa y sin tasas reducidas, con un umbral de registro lo bastante alto como para
que la microempresa quede fuera. Los alimentos \emph{no} van a tasa cero: la tasa cero
es regresiva por peso de recaudación sacrificada ---los ricos compran más comida que
los pobres en términos absolutos--- y la reparación distributiva corresponde al lado
del gasto, por la cláusula~2 del Diseño~\ref{d:benefit}, donde puede apuntarse.

\rl{3}{Impuesto sobre la renta, retenido y sencillo} Progresivo, pocos tramos, un
umbral exento alto para que la mayoría de los trabajadores no declare nada, y
retención en la fuente dondequiera que haya una fuente. El objetivo de la primera
década es cobertura y hábito, no recaudación.

\rl{4}{Impuesto sobre la propiedad} Un impuesto anual sobre la propiedad residencial y
comercial, evaluado por superficie y ubicación mediante fórmula y no por juicio de
tasación, recaudado localmente y retenido localmente según el Diseño~\ref{d:municipal}.
Es el impuesto más fácil de administrar, el más difícil de evadir y el ingreso local
natural. Es además el instrumento que hace que el regalo de la Reforma Urbana de 1960
pague los servicios que lo rodean: Cuba tiene propiedad de vivienda casi universal y un
parque habitacional que se cae, y los dos hechos son el mismo hecho.

\rl{5}{Régimen presuntivo para la microempresa} Un único pago pequeño basado en la
facturación, que sustituye a la vez al impuesto sobre la renta, al IVA y a las
contribuciones sociales, según el modelo latinoamericano del monotributo, con una vía
publicada de graduación al régimen ordinario. Si esto se hace mal el sector privado se
queda informal para siempre, que es el fracaso sureuropeo contra el que está diseñado
este capítulo.

\rl{6}{Las bases móviles se gravan poco y sencillo} Los ingresos por trabajo remoto y
de nómadas del capítulo~\ref{ch:talent} pagan una tasa plana baja sin deducciones. Una
tasa alta sobre una base que puede irse no recauda nada; una tasa baja sobre una base
que se queda recauda algo y vale la pena cumplirla.

\rl{7}{Rentas donde existan} Un gravamen al alojamiento por noche de visitante,
ingresos por arrendamiento y concesión de tierra y costa estatales, y regalías del
níquel. Véase el Diseño~\ref{d:rentrule} sobre lo que puede hacerse con lo recaudado,
que es la parte que importa.

\rl{8}{Lo que no se hace} Nada de impuesto al patrimonio ---inadministrable aquí y
expulsa exactamente el capital de la diáspora que el capítulo~\ref{ch:capital} intenta
atraer---. Ningún impuesto de salida a los emigrantes, que es un impuesto sobre el
agravio que este libro trata de saldar. Ninguna tasa nominal alta al beneficio
empresarial con exenciones negociadas, que es una máquina de corrupción con un impuesto
adosado. Ninguna afectación de ingresos a fines específicos más allá del
Diseño~\ref{d:floor}, porque fragmenta un presupuesto y después lo ata en el año
equivocado.
\end{rules}
\end{design}

\section{La trampa de la renta}

Las Partes I y II muestran el mismo fracaso tres veces: la cuota azucarera
norteamericana, los protocolos soviéticos, el petróleo venezolano. Hay un cuarto
disponible ---un auge turístico, un maná de la normalización, un pico del precio del
níquel--- y si llega se gastará en compromisos sociales corrientes, porque es lo que
ha hecho todo gobierno en esta situación, incluido este.

El Estado de bienestar tiene que financiarse con una base doméstica amplia
\emph{porque} uno financiado con renta ya ha muerto dos veces en memoria viva. No es
una preferencia por la tributación en términos ideológicos. Es la lección institucional
específica de los capítulos~\ref{ch:sovieteconomy} y \ref{ch:venezuela}.

\begin{design}{La regla de la renta}
\label{d:rentrule}
\begin{rules}
\rl{1}{Separación} Los ingresos por concesiones, regalías, privatizaciones, alivio de
deuda y cualquier maná se ingresan en un fondo de estabilización e inversión y no
están disponibles para el presupuesto corriente.

\rl{2}{Usos permitidos} El fondo financia gasto de capital, la transición de pensiones
del Diseño~\ref{d:pension}, y el apoyo contracíclico tras un huracán o un choque
externo. No financia salarios, prestaciones ni subsidios.

\rl{3}{Ancla} El presupuesto corriente se equilibra a lo largo del ciclo únicamente con
ingresos que no sean renta, con la regla de déficit publicada, vigilada por un consejo
fiscal independiente y con informe anual a la asamblea.

\rl{4}{Escape} La regla solo puede suspenderse por desastre natural declarado o choque
externo, por mayoría cualificada, con caducidad automática y una senda de retorno
publicada. Una regla fiscal con una válvula de escape discrecional no es una regla, y
una regla fiscal sin ninguna válvula de escape no sobrevive a un huracán de categoría
5 ---que Cuba va a tener---.
\end{rules}
\end{design}

\section{La aritmética de las pensiones}

Es la cifra más difícil del libro y el capítulo no debe fingir otra cosa.

La fecundidad cubana está por debajo del reemplazo desde 1978, ronda hoy 1,2 a 1,4, y
no hay ninguna cohorte en camino. Más de una cuarta parte de la población tiene más de
sesenta años. Y entre uno y dos millones de personas, desproporcionadamente en edad de
trabajar, se fueron entre 2021 y 2025. Una pensión de reparto con tasas de reemplazo
europeas no es financiable con esa razón de dependencia a ninguna tasa impositiva
plausible, y ningún diseño cambia eso.

Lo que el diseño sí puede hacer es decidir quién carga con la brecha y volver visible la
aritmética por adelantado en vez de en una crisis.

\begin{design}{Pensiones}
\label{d:pension}
\begin{rules}
\rl{1}{Dos tramos} La pensión plana universal del Diseño~\ref{d:benefit}, financiada
con impuestos, fijada para prevenir la pobreza y actualizada por fórmula; más un tramo
ligado a ingresos en cuentas individuales. El tramo plano es lo que protege a las
personas que el capítulo~\ref{ch:premises} identificó como expuestas. El tramo de
ingresos es lo que evita que el tramo plano tenga que ser generoso.

\rl{2}{Edad de retiro indexada} La edad de pensión se indexa a la esperanza de vida
mediante una fórmula legislada una sola vez y aplicada automáticamente después.
Legislar la fórmula es políticamente posible; legislar cada aumento por separado no lo
es, que es la razón por la que los países que eligieron la segunda vía no tienen los
aumentos.

\rl{3}{La inmigración es política de pensiones} El instrumento más barato para una
razón de dependencia es más gente en edad de trabajar. La migración de retorno y la
inmigración se tratan por tanto en el capítulo~\ref{ch:talent} como instrumentos
fiscales y no como cuestión cultural, y el régimen de residencia de ese capítulo está
diseñado teniendo a la vista la aritmética de este.

\rl{4}{Portabilidad} Convenios bilaterales de seguridad social con los países donde
está la diáspora ---Estados Unidos, España, México, Italia--- de modo que las
contribuciones hechas en el exterior cuenten, los retornados no sean penalizados, y los
cubanos en el exterior puedan contribuir voluntariamente. Esto convierte una diáspora
en una base de cotizantes, cosa que ningún otro instrumento hace.

\rl{5}{Contabilidad honesta} Un informe actuarial independiente a la asamblea cada año,
publicado, que declare la brecha proyectada y qué la cierra. Parte de esta brecha no se
cierra, y una cifra publicada es la manera en que un país decide qué parte cubrirá con
retiro más tardío, cuál con tasas de reemplazo menores, cuál con impuestos más altos y
cuál con migración. El déficit de pensiones cubano hoy no se publica ni se proyecta.
\end{rules}
\end{design}

\section{El personal}

Un Estado de bienestar es gente, y aquí es donde el arreglo actual está fallando más
rápido. Los médicos, enfermeros y maestros cubanos se fueron al sector turístico en los
años noventa y se están yendo del país en los años veinte. Ningún plan de financiamiento
funciona si el personal se ha ido, y ningún plan que trate la dotación como gratuita ha
sido costeado.

\begin{design}{La paga pública}
\label{d:pay}
\begin{rules}
\rl{1}{Referencia} La paga profesional pública se fija contra el salario privado y
turístico doméstico, se publica, y la revisa anualmente un organismo independiente. No
se fija contra una escala histórica en pesos, que es como un cirujano llegó a ganar
menos que un botones y como toda la sociedad aprendió que su sistema educativo era una
broma.

\rl{2}{Consistencia} Los capítulos~\ref{ch:tourism} y \ref{ch:life} quedan obligados por
esta cláusula. Una estrategia turística que eleve los salarios de servicios sin un
acuerdo correspondiente de paga pública es una estrategia para vaciar los hospitales, y
Cuba ya corrió ese experimento.

\rl{3}{Costeada} La masa salarial pública bajo la cláusula~1 es la mayor partida
individual del costo de este capítulo y se declara como tal en la tabla de
financiamiento del capítulo~\ref{ch:sequence} en vez de tratarse como un detalle
administrativo.

\rl{4}{Retener antes que reclutar} En la primera fase, el gasto va a retener al personal
existente antes que a formar sustitutos, porque una enfermera formada que se va tarda
cuatro años en reemplazarse y la salida se acumula.
\end{rules}
\end{design}

\section{La prestación, y dónde cae la línea}

La cobertura universal no exige la prestación monopólica por el Estado, y fingir que sí
le cuesta a Cuba una capacidad que no puede permitirse perder. Pero un sistema de dos
niveles en el lugar equivocado destruye aquello que se protege, porque una vez que las
clases profesional y política se van del sistema público este deja de defenderse en el
presupuesto. La línea hay que trazarla explícitamente en vez de dejarla a la práctica.

\begin{design}{Prestación pública y privada}
\label{d:delivery}
\begin{rules}
\rl{1}{Prestación privada permitida} En odontología, farmacia, óptica, procedimientos
electivos y estéticos, cuidado de mayores, rehabilitación, y educación terciaria y
vocacional no obligatoria. En ellas, prestadores privados y cooperativos operan junto al
sistema público, están licenciados e inspeccionados, y pueden contratar con él.

\rl{2}{Prestación pública protegida} La atención primaria, la atención de urgencias, la
maternidad, la salud pública y la inmunización, y la escolarización obligatoria siguen
siendo universales, públicas y gratuitas en el punto de uso, sin que se permita un
nivel privado paralelo que las sustituya. Son los servicios cuya universalidad produce
los resultados del capítulo~\ref{ch:ledgerch}, y son los primeros que un arreglo de dos
niveles vacía.

\rl{3}{Sin adelantar la cola} Ninguna instalación pública vende acceso prioritario. El
turismo de salud del capítulo~\ref{ch:life} opera en instalaciones designadas con
personal acotado, y sus ingresos van al sistema público ---y si ese acotamiento no puede
mantenerse, el programa se detiene, porque el capítulo~\ref{ch:premises} estableció que
el personal es la escasez vinculante---.

\rl{4}{Cooperativas primero} Donde se abra la prestación, las formas cooperativas y de
sociedad profesional son la opción por defecto y reciben la preferencia en la licencia,
porque distribuyen el excedente entre los profesionales y son la forma con menor
propensión a concentrarse.
\end{rules}
\end{design}

\section{Que se sostenga}

Todo lo anterior es política pública hasta que quede anclado, y el
capítulo~\ref{ch:polity} estableció qué les ocurre a las políticas públicas cubanas.

\begin{design}{El piso}
\label{d:floor}
\begin{rules}
\rl{1}{Mínimo constitucional} Un piso al gasto conjunto en salud y educación como
proporción del gasto público general, fijado en la constitución, por debajo del cual no
puede presentarse un presupuesto. Es la única afectación que permite el
Diseño~\ref{d:tax}.8, y se permite porque es el dispositivo de compromiso sobre el que
descansa el capítulo entero.

\rl{2}{Cuentas separadas} Las cuentas del seguro social se llevan separadas del
presupuesto general, se publican íntegras, y no puede tomarse prestado de ellas.

\rl{3}{Informe actuarial anual} Según el Diseño~\ref{d:pension}.5, a la asamblea,
publicado, independiente.

\rl{4}{Justiciable} Los derechos a la atención de salud y a la escolarización obligatoria
son exigibles por un individuo ante el tribunal administrativo del
Diseño~\ref{d:courts}. No el nivel de prestación ---los tribunales no deben fijar
presupuestos--- sino el hecho del acceso, y su no discriminación.

\rl{5}{Resultados publicados} Los resultados de salud y educación se publican anualmente
por provincia y, según el capítulo~\ref{ch:premises}, desagregados por raza. Cuba
declaró resuelta la cuestión racial en 1962 y después la volvió inmedible durante
treinta años, y entró en su apertura de mercado sin idea de dónde estaba parada. El
primer requisito de este diseño en ese dominio es publicar los datos.
\end{rules}
\end{design}

\section{Lo que este capítulo le cuesta al resto del libro}

Enunciado con llaneza, porque un diseño que esconde sus disyuntivas es un folleto.

Un Estado de bienestar europeo significa una carga tributaria muy por encima de la norma
latinoamericana. Cuba estará pujando por inversión contra países de la misma región con
tasas más bajas, protección laboral más débil y ningún gasto social comparable, y perderá
parte de esa puja. El capítulo~\ref{ch:capital} no puede por tanto competir en tasas y
tiene que competir en las otras cosas que los inversionistas valoran: certidumbre,
rapidez de aprobación, contratos exigibles, una aduana honesta, un tribunal en el que se
pueda perder limpiamente, y una fuerza de trabajo que sepa leer. Son más difíciles de
construir que una exoneración fiscal y son además, a diferencia de una exoneración
fiscal, valiosas por sí mismas.

Significa además que los capítulos de crecimiento tienen que funcionar. Un Estado de
bienestar de este nivel exige una base tributaria que hoy no existe, es decir, exige que
los capítulos~\ref{ch:energy} a \ref{ch:capital} tengan éxito. Este capítulo no es una
alternativa a la estrategia de crecimiento. Es la razón de ella.
